\documentclass[../main.tex]{subfiles}  % 使用主文件的文档类

\begin{document}

\section{Symbolic Modeling of Entropy-Based Security}


% 在上一章节中，我们回顾了先进的验证工具，并详细介绍了本文使用的验证工具Tamarin prover.
% 在本章节中，我们将使用Tamarin prover对协议中的低熵密钥以及广义的暴力破解行为进行建模。

% 在这一章节，我们对密钥空间受限的密钥的安全性和弱点进行了系统的概括，这是第四章符号模型的基础。
% 在前一章节中，我们回顾了文献中低熵密钥造成的安全问题。然而，对密钥熵的分析不存在现有的任何符号建模的工作中。通常在符号模型中，密码是完美的，攻击者无法从密文中获得任何信息。对于我们的方法，我们首先确定了低熵密钥对协议造成安全隐患的核心属性：保密性。我们回顾的攻击中，攻击者的目的就是通过获取密钥来破坏协议。因此，我们将密钥的保密性作为核心关注点，对其进行研究。

In the previous section, we reviewed advanced verification tools and introduced the verification tool Tamarin prover used in this article in detail.
In this section, we will use Tamarin prover to model low entropy keys in the protocol and generalized brute force cracking behavior.

\subsection{Entropy and Brute Force Attack}

% 熵是信息论中的一个重要概念，由香农在1948年提出，用于衡量系统的随机性和不确定性。在密码学中，密钥的熵是衡量密钥强度的重要指标，一个在密钥空间均匀分布的密钥的熵通常可以用以下的公式表示：
% H = log_2(N)
% 其中N是密钥空间的大小。该公式说明了，均匀分布的密钥的熵就是该密钥的比特数。那么对于想要通过穷尽搜索的方式找到密钥的攻击者来说，他需要尝试的密钥的数量就是2^H，其中H是密钥的熵。因此，密钥的熵越高，攻击者找到密钥的难度就越大。


Entropy is a fundamental concept in information theory, introduced by Shannon in 1948 {cite}, to quantify the randomness and uncertainty within a system.
In the domain of cryptography, the entropy of a key serves as a crucial metric for assessing its strength.
 Formally, consider a uniformly distributed key $K$ within a key space of size $N$.
The entropy of such a key can be represented by \autoref{eq:entropy_definition}.
\begin{equation} 
    H(K) = \log_2(N) \label{eq:entropy_definition}
\end{equation}
where $H(K)$ denotes the entropy of the key $K$.
This equation illustrates that the entropy of a uniformly distributed key is equivalent to its bit length.
For an adversary employing exhaustive search techniques to locate the key, the number of keys that must be tested is $2^{H(K)}$.
Therefore, an increase in the key's entropy directly correlates with a higher level of difficulty for an attacker attempting to restore the key.

% 这种遍历密钥空间的方法，不断地尝试密钥，直到找到正确的密钥，就是暴力破解攻击。通过是否需要与目标系统进行实时交互可以分为在线攻击与离线攻击。在线的攻击，攻击者需要不断地将尝试的密钥发送给目标系统以判断密钥是否正确。离线攻击，通常是通过获取目标系统的密文，攻击者可以在不与目标系统交互的情况下，尝试所有的密钥通过解密密文的方式来鉴别密钥。通常来说在线的暴力破解攻击主要受限于网络速度和目标系统的响应时间，而离线的攻击主要受限于攻击者的计算速度。离线暴力破解攻击的运行效率完全掌握在攻击者手中，密钥因此需要具有足够的熵来抵御这种攻击。因此本文主要研究的就是离线的暴力破解攻击，后文提到的暴力破解攻击如果没有特别制定，均指离线的暴力破解攻击。

The method of exhaustively searching the key space by continually attempting different keys until the correct one is identified is termed a brute-force attack.
Brute-force attacks can be classified into online and offline categories based on the necessity for real-time interaction with the target system.
In an online attack, the adversary must persistently transmit the attempted keys to the target system to verify their correctness.
Conversely, in an offline attack, the adversary typically acquires the ciphertext from the target system and can attempt all possible keys by decrypting the ciphertext, thereby obviating the need for interaction with the target system.

% \begin{figure}[htbp]
%     \centering
%     
\includegraphics[width=\columnwidth]{img/online-brute-force} 
%     \caption{online-brute-force}
%     \label{fig:online-brute-force}
% \end{figure}

Typically, the efficacy of online brute-force attacks is constrained by network latency and the response time of the target system.
In contrast, offline attacks are predominantly limited by the computational power available to the attacker.
The operational efficiency of offline brute-force attacks is entirely governed by the attacker's capabilities; therefore, the cryptographic key must possess adequate entropy to resist such attacks effectively Consequently, this paper primarily focuses on the investigation of offline brute-force attacks.
Throughout the subsequent sections, any reference to brute-force attacks, unless explicitly stated otherwise, pertains to offline brute-force attacks.

% \begin{figure}[htbp]
%     \centering
%     
\includegraphics[width=\columnwidth]{img/offline-brute-force} 
%     \caption{offline-brute-force}
%     \label{fig:offline-brute-force}
% \end{figure}


\subsection{Key Breaking Oracles}

% 现有的符号模型采用完美密码假设，攻击者无法从密文中获取任何信息。然而Cas{cite}, Min{cite}，xxx{cite}等人在建模特定协议时，通过一种特殊的规则，模拟了攻击者访问Oracle的行为，向攻击者某些密钥，一个简单的规则如下：
% 使用这样的规则，可以赋予攻击者破解某些密钥的能力。然而，对于低熵密钥而言，这种方法不具有普适性，因为它需要手动地分析哪些低熵密钥是可以被破解的。

Existing symbolic models typically adopt the perfect cipher assumption, which posits that an attacker cannot derive any information from the ciphertext.
However, researchers such as Cas{cite}, Min{cite}, and xxx{cite} have advanced the modeling of certain protocols by introducing a special rule.
This rule simulates the behavior of attackers accessing an oracle, thereby allowing them to obtain certain keys.
An example of such a rule is as follows:
\begin{flalign*}
    \text{\ruleNameText{Orc}} & \text{\ruleNameText{ale}} :                                        &\\
    \TFacts{ &\text{\TFact{f}}(s, pk1, pk2), \text{\TFact{In}}(pk1), \text{\TFact{In}}(pk2)   } \TEmpActFact{ } \\
    \TFacts{ &\text{\TFact{Out}}(s)}    
\end{flalign*}
Employing these rules can enable attackers to compromise specific keys.
Nonetheless, this approach is not universally applicable to low-entropy keys, as it necessitates a manual evaluation to identify which low-entropy keys are susceptible to being compromised.

% 为了通用化和自动化，我们需要一个统一的方法，自动分析出协议中所有可被破解的低熵密钥。类似的，我们假设攻击者可以访问一个密钥破解的oracle，来建模攻击者的穷尽搜索能力，这个oracle如下：

To achieve generalization and automation, it is essential to develop a unified method capable of automatically analyzing all low-entropy keys within the protocol that are susceptible to compromise.
In a similar vein, we assume the attacker has access to a key-breaking oracle, which serves to model the attacker's exhaustive search capabilities.
The oracle is shown as follows:
\begin{figure}[!h]
    \centering
    \includegraphics[width=0.7755\columnwidth]{img/general-oracle} 
\end{figure}

% 在这个oracle中，攻击者可以提供一个有低熵密钥参与的函数，如果能提供除密钥外的其他参数，则攻击者就可以获得该秘密的低熵密钥。直接建模这样的Oracle是困难的，一方面，Tamarin 的 pattern matching 是精确的，我们无法使用一个 function symbol 匹配所有密钥参与的算法，另一方面，攻击者需要输入除密钥外所有参与运算的参数，Oracle动态地根据具体算法来决定需要的参数。我们的方法设计了两个更加基础的Oracle，共同完成密钥破解的过程，一个是算法破解Oracle，另一个是值破解Oracle。
\noindent In this oracle, an attacker can provide a function that incorporates a low-entropy key.
If all parameters involved in the computation except for the key are supplied, the attacker can retrieve the secret low-entropy key.
Directly modeling such an oracle presents significant challenges.
Firstly, Tamarin's pattern matching is exact, which precludes the use of a single function symbol to match all algorithms involving keys.
Secondly, the attacker must input all parameters involved in the computation aside from the key, and the oracle dynamically determines the necessary parameters based on the specific algorithm in question.
Our methodology introduces two more fundamental oracles to collaboratively facilitate the key-breaking process.


% a value-breaking oracle and an algorithm-breaking oracle.
\subsubsection{Value-breaking Oracle}
% 数值破解Oracle可以有限地赋予攻击者破解低熵密钥的能力，前提是攻击者可以找到低熵密钥作为唯一变量的计算结果。该Oracle的流程如图\autoref{label}所示。攻击者可以向Oracle提供密钥作为唯一变量的计算结果，如果该密钥的熵低，那么Oracle会向攻击者提供该低熵密钥的数值。
The value-breaking oracle grants attackers the limited capability to break low-entropy keys, assuming the attacker can identify a computation result in which the low-entropy key is the only variable.
The procedure for this oracle is illustrated in Figure \autoref{fig:value-oracle}.
Attackers can submit the computation result, where the key is the sole variable, to the oracle.
If the key exhibits low entropy, the oracle will provide the attacker with the value of this low-entropy key.

 \begin{figure}[h]
    \centering
    \includegraphics[width=0.4653\columnwidth]{img/oracle1} 
\end{figure}

\subsubsection{Algorithm-breaking Oracle}
% 图4展示了算法破解oracle的执行流程，攻击者可以通过提供一个有密钥参与的函数，如果能提供除密钥外的其他参数，算法破解Oracle可以找到另一个函数，满足Eq.
\ref{eq:entropy_definition}.
\autoref{fig:algorithm-oracle} delineates the execution process of the algorithm-breaking oracle.
An attacker is capable of presenting a function $f$ that incorporates a key $k$.
Given that the attacker can provide all parameters $x_1, x_2, ...
, x_n$ except for  $k$,  the oracle can identify another function $f'$ that satisfies  \autoref{eq:function_cracking}.
\begin{equation} 
    f'(k) = f(\langle k, x_1, x_2, ...
, x_n \rangle) \label{eq:function_cracking}
\end{equation}
% 当攻击者使用这个oracle，确定了使用低熵密钥作为唯一变量的函数，并获取了对应的计算结果，那么攻击者就可以使用值破解Oracle来获取低熵密钥的数值。同样为了通用性，输入函数是不确定的，我们无法直接建模这个Oracle，我们将在下一小节介绍如何通过一系列规则来建模algorithm-breaking oracle。
When an attacker utilizes this oracle to identify a function that employs a low-entropy key as the sole variable and acquires the corresponding computation result, the attacker can subsequently utilize the value-breaking oracle to obtain the value of the low-entropy key.
Similarly, due to the uncertainty of the input function, it is infeasible to directly model this oracle.
In the subsequent subsection, we will elucidate the method to model the algorithm-breaking oracle through a set of message derivation rules.
\begin{figure}[htbp]
    \centering
    \includegraphics[width=0.99\columnwidth]{img/oracle2} 
\end{figure}


\subsection{Oracle Modeling} \label{sec:oracle-modeling}

% 我们通过一系列规则来建模算法破解Oracle。为抽象化oracle内部消息的派生过程，我们提出了消息派生事实的概念，以代替具体的函数符号。此外，为了建模algorithm-breaking oracle，我们还需要引入消息派生规则，帮助oracle推导出攻击者视角下所有可见的消息派生事实。在此基础之上，我们将介绍如何建模value-breaking oracle。 最后，我们将讨论消息派生规则执行的一些限制。


We model the previous oracles through a series of rules.
To abstract the process of message derivation within the Oracle, we introduce the concept of message derivation facts, replacing specific function symbols.
Furthermore, to effectively model the algorithm-breaking Oracle, it is necessary to incorporate message derivation rules that enable the Oracle to deduce all message derivation facts visible from the attacker's perspective.
Building on this foundation, we will present our approach to modeling the value-breaking Oracle.

% Finally, we will discuss certain restrictions associated with the execution of message derivation rules.

% 消息派生事实 在algorithm-breaking oracle的形式化中，我们将密钥的计算抽象为派生关系，使用消息派生fact来表示，每个派生Fact都是一个消息到一个消息的映射。具体来说，对于一个单参数函数$f$ 以及变量 $x$，我们可以将$f$抽象为如下的Fact：
% \[
%     MessageDrive(x, f(x))
% \]
% % 为了形式上的统一，对于多变量的函数，建模为
% \[
%     MDerive(<x_1, x_2, ..., x_n>,f(<x_1, x_2, ..., x_n>))
% \]
% ，其中<>是一个n个变量的function，是n个函数变量的拼接。类似的，复合函数$f \circ  g$，可以被建模为如下的消息派生事实：
% \[
% MDerive(x, f(g(x)))
% \]
% 
% algorithm-breaking oracle维护了所有函数的集合。只要明确了两个消息间存在派生关系，oracle可以直接找到具体的映射函数。 在此基础之上，如果oracle需要找到某个密钥$k$作为唯一变量的函数，它所要做的只要找到如下fact。
% \[
% MDerive(k, m)
% \]

\subsubsection{Message Derivation Fact} In the formalization of an algorithm-breaking oracle, key computation is abstracted into derivation relationships, represented by message derivation facts.
Each derivation fact maps one message to another.
Specifically, for a one-argument function $f$ and an argument $x$, $f$ can be abstracted as the following fact:
\[
    \text{\TFact{MDerive}}(x, f(x))
\]
For the sake of consistency, multi-argument functions are modeled as:
\[
    \text{\TFact{MDerive}}(\langle x_1, x_2, ..., x_n \rangle, f(\langle x_1, x_2, ..., x_n \rangle))
\]
where $\langle \cdot \rangle$ denotes a function of $n$ arguments, representing the concatenation of $n$ function arguments.
Similarly, a composite function $f \circ g$ can be modeled as the following message derivation fact:
\[
\text{\TFact{MDerive}}(x, f(g(x)))
\]
The algorithm-breaking oracle maintains a collection of all functions.
Once a derivation relationship between two messages is identified, the oracle can directly locate the specific mapping function.
Based on this framework, if the oracle needs to identify a function where a certain key $k$ is the sole argument, it suffices to find the following fact:
\[
\text{\TFact{MDerive}}(k, m)
\]
where $m$ is a message derived from $k$.

% 消息派生规则 给定oracle $O$和敌手$A$，$df_O$代表oracle拥有的所有消息派生关系事实的集合，$dm_A$代表敌手所有可推断消息的集合。 在敌手$A$ 开始访问 $O$ 时，$O$ 会获取协议当前执行过程中所有的消息派生事实$\mathcal{P}(MessageDeriveFact)$，初始化$df_O$。$df_O$被如下规则被inductively定义：
% \subsubsection{Message Derivation Rule}  Consider an oracle $ \mathcal{O} $ and an adversary $ \mathcal{A} $.
The set $ df_\mathcal{O} $ denotes all message derivation facts known to the oracle, while $ dm_\mathcal{A} $ represents the set of all deducible messages that the adversary can know.
When the adversary $ \mathcal{A} $ initiates interaction with the oracle $ \mathcal{O} $, the oracle gathers all message derivation facts from the current execution of the protocol $\mathcal{P}$, denoted by $df_\mathcal{P}$, and uses these to initialize $ df_\mathcal{O} $.
The set $ df_\mathcal{O} $ is inductively defined by the following rules:

% \[
% \infer
%     {\text{\TFact{MDerive}}(m_1, m_2) \in df_\mathcal{O}}
%     {\text{\TFact{MDerive}}(m_1, m_2) \in df_\mathcal{P}}
% \]

% \[
% \infer
%     {\text{\TFact{MDerive}}(m_1, m_3) \in df_\mathcal{O}}
%     {\text{\TFact{MDerive}}(m_1, m_2) \in df_\mathcal{O} & \text{\TFact{MDerive}}(m_2, m_3) \in df_\mathcal{O}}
% \]

% \[
% \infer
%     {\text{\TFact{MDerive}}(m_2, m_3) \in df_\mathcal{O}}
%     {\begin{array}{c} 
%         m_1 \in dm_\mathcal{A} \\ 
%         \text{\TFact{MDerive}}(\langle m_1, m_2 \rangle, m_3) \in df_\mathcal{O}
%     \end{array}}
% \]

% \[
% \infer
%     {\text{\TFact{MDerive}}(m_1, m_3) \in df_\mathcal{O}}
%     {\begin{array}{c} 
%         m_2 \in dm_\mathcal{A} \\ 
%         \text{\TFact{MDerive}}(\langle m_1, m_2 \rangle, m_3) \in df_\mathcal{O}
%     \end{array}}
% \]

% 消息派生规则 在攻击者开始访问oracle时，oracle 会获取协议当前执行过程中的所有派生事实，并使用如下的规则推断出所有对攻击者而言可计算的消息派生事实：
% 其中，攻击者需要不断提供所 有已知的term，以辅助oracle推断出所有可见的派生事实。
\subsubsection{Message Derivation Rule}  
Upon the initiation of the attacker's access to the oracle, the oracle acquires all derivation facts present in the current execution process of the protocol.
It then infers all existing derivation facts that are computable from the attacker's perspective using the following rules:
\begin{flalign*}
    \text{\ruleNameText{Orc}} & \text{\ruleNameText{aleMsgDerive1}} :                               &\\
    \TFacts{ &\text{\TFact{MDerive}}(m_1, m_2), \text{\TFact{MDerive}}(m_2, m_3) } \TEmpActFact{ }  \\
    \TFacts{ &\text{\TFact{MDerive}}(m_1, m_3) }        
    \\
    \text{\ruleNameText{Orc}} & \text{\ruleNameText{aleMsgDerive2}} :                               &\\
    \TFacts{ &\text{\TFact{MDerive}}(\langle m_1, m_2 \rangle, m_3), \text{\TFact{In}}(m_1) } \TEmpActFact{ }  \\
    \TFacts{ &\text{\TFact{MDerive}}(m_2, m_3) }
    \\
    \text{\ruleNameText{Orc}} & \text{\ruleNameText{aleMsgDerive3}} :                               &\\
    \TFacts{ &\text{\TFact{MDerive}}(\langle m_1, m_2 \rangle, m_3), \text{\TFact{In}}(m_2) } \TEmpActFact{ }  \\
    \TFacts{ &\text{\TFact{MDerive}}(m_1, m_3) }
\end{flalign*}
In these rules, the attacker is required to continuously supply all known terms to assist the oracle in inferring all visible derivation facts.

\subsubsection{Value Breaking Oracle} 
% oracle 破解密钥的最后一步，检查目标的key是否是低熵，且攻击者掌握其作为唯一变量的函数输出，如果是，oracle会返回该密钥的数值。这个规则如下：
% 其中，k是攻击者的目标密钥，m是密钥的计算结果，$\text{\TFact{LowEntropy}}(k)$表示k是低熵密钥，在建模具体协议时被手动指定。
In the concluding step of the oracle's key-breaking process, the oracle verifies whether the target key exhibits low entropy and whether the attacker has knowledge of a function output in which the key is the sole variable.
If these conditions are satisfied, the oracle will return the value of the key.
This rule can be stated as follows:
\begin{flalign*}
    \text{\ruleNameText{Orc}} & \text{\ruleNameText{aleKeyCrack}} :                               &\\
    \TFacts{ &\text{\TFact{MDerive}}(k, m), \text{\TFact{LowEntropy}}(k), \text{\TFact{In}}(m) } \TEmpActFact{ }  \\
    \TFacts{ &\text{\TFact{Out}}(k) }        
\end{flalign*}
where $k$ denotes the attacker's target key, $m$ represents the computation result from the key, and $\text{\TFact{LowEntropy}}(k)$ labels that $k$ is a low-entropy key, which is manually specified during modeling a certain protocol.
Such a fact is essential for the oracle to determine whether the key's entropy is low.



\end{document}
